\section{Short proofs and declared reductions}
\label{app:short-proofs}
\subsection{Stochastic crowning}
\label{app:stochastic-crowning}

\subsection{Scope}
\label{app:stoch-scope}
A\(\Omega\)D stochastic accounting is count-normalized boundary accounting over a declared field, boundary, window, and fractal address.


\aodliteraturenote{Finite Markov kernels: Norris~\cite{norrisMarkov}, Levin--Peres--Wilmer~\cite{levinPeresWilmer}.}

\subsection{Setup}
\label{app:stoch-setup}
Declare a finite history family on the field and window:
\[
\mathcal H_{\mathcal F}(\omega).
\]
Each history has the form
\[
h=(h_{\mathrm{past}},h_{\mathrm{present}},h_{\mathrm{future}}).
\]

\subsection{Statement(s)}
\label{app:stoch-statements}

\subsubsection{History family}
\label{app:stoch-history-family}
The family \(\mathcal H_{\mathcal F}(\omega)\) is the declared finite support for stochastic crowning on the window.

\subsubsection{Crowning}
\label{app:stoch-crowning}
The crowning map selects the present:
\[
\operatorname{Crown}(h)=h_{\mathrm{present}}.
\]

\subsubsection{Alpha advantage}
\label{app:stoch-alpha-advantage}
At one declared observation restriction, partition histories into mutually
exclusive statuses: crowned-alpha and not yet returned, completed Omega
return, or unresolved. A history is counted in exactly one current status;
its earlier alpha event remains in provenance but is not counted again as a
second current history. With this partition, define
\[
N_\alpha=
\#\{h:\ h\text{ has current alpha status}\},
\]
\[
N_\Omega=
\#\{h:\ h\text{ has completed-return status}\},
\]
\[
N_0=
\#\{h:\ h\text{ remains unresolved}\}.
\]
The total count is
\[
N_{\mathrm{tot}}=N_\alpha+N_\Omega+N_0,
\qquad
N_{\mathrm{tot}}>0.
\]
The alpha advantage is
\begin{equation}
\Delta_\alpha(\omega)=
\frac{N_\alpha-N_\Omega}{N_{\mathrm{tot}}}.
\label{app:stochastic:eq:alpha-advantage}
\end{equation}

\subsubsection{Retained-history family}
\label{app:stoch-retained-history-family}
For channel \(e\), the retained-window indicator is
\[
\mathbf 1_{\mathrm{ret}}(e,\omega)=
\begin{cases}
1, & \mathrm{TTL}_e(\omega)\ge |\omega|_e,\\
0, & \mathrm{TTL}_e(\omega)< |\omega|_e.
\end{cases}
\]
The history retention indicator is
\[
\mathbf 1_{\mathrm{ret}}(h,\omega)=1
\]
when every retained channel required by \(h\) satisfies the declared TTL policy on \(\omega\). The same mutually exclusive status partition is used after filtering. The retained-history family is
\[
\mathcal H_{\mathcal F}^{\mathrm{ret}}(\omega)
=
\{h\in\mathcal H_{\mathcal F}(\omega):
\mathbf 1_{\mathrm{ret}}(h,\omega)=1\}.
\]
Define
\[
N_{\alpha}^{\mathrm{ret}}
=
\#\{h\in\mathcal H_{\mathcal F}^{\mathrm{ret}}(\omega):
h\text{ has current alpha status}\},
\]
\[
N_{\Omega}^{\mathrm{ret}}
=
\#\{h\in\mathcal H_{\mathcal F}^{\mathrm{ret}}(\omega):
h\text{ has completed-return status}\},
\]
\[
N_{\mathrm{tot}}^{\mathrm{ret}}
=
|\mathcal H_{\mathcal F}^{\mathrm{ret}}(\omega)|.
\]
The retained alpha advantage is defined when \(N_{\mathrm{tot}}^{\mathrm{ret}}>0\):
\begin{equation}
\Delta_{\alpha}^{\mathrm{ret}}(\omega)
=
\frac{N_{\alpha}^{\mathrm{ret}}-N_{\Omega}^{\mathrm{ret}}}{N_{\mathrm{tot}}^{\mathrm{ret}}}.
\label{app:stochastic:eq:ret-alpha-advantage}
\end{equation}

\subsection{Derivation}
\label{app:stoch-derivation}
For a temporal cycle group \(W\), use the window-clipped duration participation \(\rho^D_{\omega,e}\) inside the temporal SADAR burden:
\[
P_W(\partial F,\omega)
=
\sum_{e\in W}
C_{3,e}
\rho^D_{\omega,e}(\partial F,\omega)
|\Delta^{\chi\mathrm{lag}}_e|.
\]
The notation \(\mathbb E_{\mathcal H}\) denotes finite arithmetic average over the declared history family \(\mathcal H_{\mathcal F}(\omega)\):
\[
\mathbb E_{\mathcal H}[P_W]
=
\frac{1}{N_{\mathrm{tot}}}
\sum_{h\in\mathcal H_{\mathcal F}(\omega)}
P_W(h).
\]

\subsection{Falsification}
\label{app:stoch-falsification}
If a declared history family has \(N_{\mathrm{tot}}=0\), the alpha-advantage expression is undefined on that scope. If TTL filtering changes the retained support, retained counts must be recomputed before using \(\Delta_{\alpha}^{\mathrm{ret}}\).

\aodprovenance{\secref{subsec:crown},
\eqnref{app:stochastic:eq:alpha-advantage},
\eqnref{app:stochastic:eq:ret-alpha-advantage},
\secref{subsec:support-enclosure-retention-condition}.}

\aodremark{The probability terms are normalized counts on the declared history family.}
\subsection{Scoped flux cancellation}
\label{app:scoped-flux-cancellation}
\label{app:symbolic-cancellation-check}

\subsection{Scope}
\label{app:scoped-flux:scope}
This App records symbolic support for finite Stokes with separately typed A\(\Omega\)D pressure-resolved flux values. The scope is a finite declared region \(R\) and boundary/window \(B=(\partial F,\omega)\).

\subsection{Statement(s)}
\label{app:scoped-flux:statements}

\[
\phi_{ij}\mapsto \phi_{ij}(B),
\qquad
\phi_{ij,\AO}(B)=p^D_{ij,\AO}(B)A_{ij,\AO}(B).
\]
The scoped boundary pairing is retained through the AFC Stokes identity:
\begin{equation}
\sum_{i\in R}\operatorname{div}(i;B)
=
\sum_{\substack{i\in R\\j\notin R}}\phi_{ij}(B).
\label{app:scoped-flux:eq:afc-pairing}
\end{equation}
The Stokes residual is
\begin{equation}
\mathcal E(R,B)
=
\sum_{i\in R}\operatorname{div}(i;B)
-
\sum_{\substack{i\in R\\j\notin R}}\phi_{ij}(B).
\label{app:scoped-flux:eq:residual}
\end{equation}
A valid declared scope has \(\mathcal E(R,B)=0\).

\subsection{Falsification}
\label{app:scoped-flux:falsification}
If the symbolic residual is nonzero, the declared scope has a missing channel, sign error, incomplete boundary, or undeclared sheddic path.

\aodprovenance{\eqnref{eq:b-scoped-flux},
\eqnref{eq:b-scoped-flux-antisymmetry},
\eqnref{eq:b-scoped-divergence},
\eqnref{eq:b-scoped-stokes},
\eqnref{eq:stokes-audit-residual},
AFC \S4.2.}

\subsection{Reflection Duration on \texorpdfstring{\(Q_4\)}{Q4}}
\label{app:combinatorics-rd-tests}

\subsection{Scope}
\label{app:rd-tests:scope}
This App records graph support, ternary fibre, Markov kernel, path histories, RD/RCD calculation, and integer support reductions used by the A\(\Omega\)-internal SADAR pressure and attention tests.

\subsection{Setup}
\label{app:rd-tests:setup}
The finite graph support is
\[
V(Q_4)=\{0,1\}^4,
\qquad
E(Q_4)=\{(x,y):d_H(x,y)=1\},
\]
with
\[
d_H(x,y)=\#\{k:x_k\ne y_k\}.
\]
A local edge state carries the ternary fibre
\[
\sigma_t(e)\in\{-1,0,+1\},
\]
where \(+1\) is out-bound, \(0\) is hinge / dwell, and \(-1\) is return-bound.

\subsection{Choice: Markov kernel}
\label{app:rd-tests:markov-kernel}
For an already admitted support with positive normalizer, the downstream kernel of \eqref{eq:admitted-kernel} satisfies
\[
P_B((e',\sigma')\mid(e,\sigma))\ge0,
\qquad
\sum_{(e',\sigma')}P_B((e',\sigma')\mid(e,\sigma))=1,
\]
with
\[
P_B((e',\sigma')\mid(e,\sigma))=0
\quad\text{outside admissible Hamming--1 continuations},
\qquad
\sigma'\in\{-1,0,+1\}.
\]
The isotropic choice is
\[
P_B^{\mathrm{iso}}((e',\sigma')\mid(e,\sigma))=
\frac{1}{N_B(e,\sigma)}
\]
over nonempty admitted, hinge-mediated successors; an empty positive-weight set has no normalized kernel. The anisotropic choice is
\[
P_B^{\mathrm{ani}}((e',\sigma')\mid(e,\sigma))=
\frac{w_B(e',\sigma')}
{\sum_{(a,\vartheta)\in\operatorname{Adm}_B(e,\sigma)}w_B(a,\vartheta)},
\]
where the weights may depend on \(\Delta\chi^{\mathrm{lag}}_{e'}(B)\), \(i_{e'}(B)\), and \(R^{\mathrm{asym}}_{e'}(B)\).

\subsection{Statement(s)}
\label{app:rd-tests:statements}
A path history is
\[
h=((e_0,\sigma_0),\ldots,(e_m,\sigma_m)).
\]
Reflection duration is first typed as branch/hinge/branch support:
\begin{equation}
RD_h(B)
=
\Pi^-_h(B)
+
H_{\mu,h}(B)
+
\Pi^+_h(B).
\label{app:rd:eq:branch-hinge-rd}
\end{equation}
Here \(\Pi^-_h\) is first-branch support, \(H_{\mu,h}\) is the hinge recurrence contribution, and \(\Pi^+_h\) is reflected return-branch support. The full-trace specialization, when the declared window contains each support occurrence exactly once, is
\[
RD_h(B)=
\sum_{t\in\omega}
\mathbf 1[h\text{ retained, pending, or returning at }t].
\]
Outside that full-trace specialization the indicator sum is window participation, not full RD. The coupling object is
\[
\operatorname{RCD}_h(B)=RD_h(B)\bowtie_B R^{\mathrm{asym}}_h(B),
\qquad
\rho^D_h(B)=\operatorname{dur}(\operatorname{RCD}_h(B)).
\]
After declaring a common exact ordered count domain for duration and window, the window-clipped scalar participation and pending overhang are
\[
\rho^D_{\omega,h}(B)=\min(\rho^D_h(B),|\omega|_h),
\]
\[
Q^{D,\mathrm{dur}}_h(B)=\max(0,\rho^D_h(B)-|\omega|_h),
\qquad
Q^D_h(B)=C_{3,h}Q^{D,\mathrm{dur}}_h(B).
\]

\subsection{Walk support corollaries}
\label{app:rd-tests:corollaries}
A\(\Omega\)D records local family support by
\begin{equation}
\Pi(p{:}q)=q^p+q.
\label{app:rd:eq:path-support}
\end{equation}
Here \(q^p\) is the memoryless \(q\)-direction walk count and \(q\) is the retained Hamming--1 edge-shell term. The branch/hinge/branch reflected-return support is
\begin{equation}
RD(p{:}q;B)=2\Pi(p{:}q)+H_\mu(B)
\label{app:rd:eq:reflected-return}
\end{equation}
when the branch supports are symmetric. The single-hinge reduction is
\begin{equation}
H_\mu(B)=1
\quad\Longrightarrow\quad
RD(p{:}q)=2\Pi(p{:}q)+1.
\label{app:rd:eq:single-hinge-return}
\end{equation}
For \(3{:}4\), under the single-hinge reduction,
\begin{equation}
\Pi(3{:}4)=4^3+4=68,
\qquad
RD(3{:}4)=2\cdot68+1=137.
\label{app:rd:eq:three-four-rd}
\end{equation}
The transition from \(4^3\) to \(4^3+4\) is A\(\Omega\)D's retained edge-shell accounting. The transition from \(\Pi\) to \(2\Pi+H_\mu\) is A\(\Omega\)D's branch/hinge/branch reflected-return closure. The single-hinge form \(2\Pi+1\) is a reduction, not the general rule.

\subsection{Integer support reductions}
\label{app:rd-tests:support-reductions}
The declared integer reductions are
\[
1{:}2{:}6,
\quad
1{:}2{:}7,
\quad
1{:}2{:}8,
\quad
1{:}2{:}9,
\quad
1{:}2{:}10,
\quad
3{:}3{:}6,
\quad
3{:}4{:}6.
\]
The following scalar reductions declare neutral reflection, compatible count windows, and their closure thresholds. The \(1{:}2{:}7\) reduction is closure-dependent:
\[
P^D(1{:}2{:}7)=3\min(7,9)=21.
\]
If closure is \(1{:}2{:}8\), then \(21-24=-3\). If closure is \(1{:}2{:}6\), then \(21-18=3\).
The \(1{:}2{:}9\) reduction is the nine-window surplus reduction:
\[
P^D=27,
\qquad
C^{\mathrm{close}}=24,
\qquad
X_{\mathrm{shedding}}=3.
\]


\aodliteraturenote{Dyck/Catalan path combinatorics: Stanley~\cite{stanleyCatalan}, Flajolet--Sedgewick~\cite{flajoletSedgewick}.\\
Markov kernels and random walks: Norris~\cite{norrisMarkov}, Levin--Peres--Wilmer~\cite{levinPeresWilmer}, Lawler--Limic~\cite{lawlerLimic}.}

\aodremark{The cited combinatorics is introduced at the first \(Q_4\) / Hamming--1 derivation in \secref{subsec:four-edge-x-witness}. This App uses that support to state A\(\Omega\)D's retained edge-shell term \(+q\) and branch/hinge/branch reflected-return closure \(\Pi^-+H_\mu+\Pi^+\). The ternary hinge fibre is carried on the finite tesseract witness; it is not identical with the tesseract support.}

\subsection{Walk-support audit}
\label{app:rd-tests:audit}
A walk-support reduction records declared first-branch support, hinge recurrence, reflected return-branch support, the retained edge-shell term, and the A\(\Omega\)D branch/hinge/branch closure. The single-hinge accessor \(2\Pi+1\) is used when the single-hinge symmetric reduction is declared. A kernel reduction records a normalized probability family over admissible Hamming--1 continuations.

\aodprovenance{\secref{subsec:four-edge-x-witness}, \eqref{eq:tesseract-admissible-boundary}, \eqnref{eq:branch-hinge-rd}, \eqnref{eq:branch-hinge-rd-symmetric}, \secref{app:rd-tests:corollaries}.}
